\documentclass[11pt,a4paper]{article}
\usepackage[utf8]{inputenc}
\usepackage[T1]{fontenc}
\usepackage[spanish]{babel}
\usepackage[landscape,margin=1.8cm]{geometry}
\usepackage{longtable,booktabs,array}
\usepackage{hyperref}

\begin{document}

\begin{center}
{\Large\bfseries FORM-TL-01: Taller de tipos de requisitos}\\[2pt]
Universidad Técnica Estatal de Quevedo \quad Ingeniería de Requerimientos (ISR-401)
\end{center}

\section{Datos del estudiante}
\begin{tabular}{@{}ll@{}}
\textbf{Estudiante:} & Sánchez Centeno Roselyn Andreina \\
\textbf{Asignatura:} & Ingeniería de Requerimientos (ISR-401), 4.º nivel, Paralelo B \\
\textbf{Docente:} & Gleiston Cicerón Guerrero Ulloa, PhD \\
\textbf{Actividad:} & FORM-TL-01, Taller de tipos de requisitos \\
\textbf{Caso:} & Plataforma de gestión de citas médicas \\
\end{tabular}

\section{Clasificación de los 20 enunciados}
% Aquí irá la tabla (Paso 3)

\section{Requisitos de calidad cuantificados}
% Aquí irán los 5 requisitos (Paso 4)

\section{Discrepancias con el compañero}
% Aquí irá la tabla (Paso 5)

\section{Declaración de uso de IA}
% Aquí irá la declaración (Paso 6)

\begin{thebibliography}{9}
\bibitem{iso29148} ISO/IEC/IEEE, \emph{ISO/IEC/IEEE 29148:2018, Systems and software engineering---Life cycle processes---Requirements engineering}, International Standard, 2018.
\bibitem{iso25010} ISO/IEC, \emph{ISO/IEC 25010:2023, Systems and software engineering---Systems and software Quality Requirements and Evaluation (SQuaRE)---Product quality model}, International Standard, 2023.
\end{thebibliography}

\end{document}
